\section{Circuits as minimal evidential redundancy}
\label{sec:matroid-circuits}

Gate~2 identified the exact topological boundary at which the coarse 4-PEL
closure also satisfies matroid exchange. Gate~3 turns from closure geometry to
combinatorial dependence. In matroid theory a circuit is a minimal dependent
set \cite{Whitney1935,Oxley2011}. For the two-channel 4-PEL projection this
notion admits a particularly sharp interpretation: a circuit is a smallest body
of evidence containing a coarse redundancy that disappears as soon as either of
its members is removed.

\begin{definition}[Channel dependence]
Let \(\chi:E\to\{+,-\}\) be the coarse evidence polarity. A body of evidence
\(A\subseteq E\) is \emph{channel-dependent} when
\[
  \exists x,y\in A\quad x\neq y
  \quad\text{and}\quad
  \chi(x)=\chi(y).
\]
It is channel-independent when no such pair exists. A \emph{channel circuit}
is a channel-dependent set with no proper channel-dependent subset.
\end{definition}

The definition is equivalent to the expected closure formulation.

\begin{proposition}[Closure witness for dependence]
For every \(A\subseteq E\),
\[
  A\text{ is channel-dependent}
  \quad\Longleftrightarrow\quad
  \exists x\in A:\
  x\in\operatorname{cl}_{\chi}(A\setminus\{x\}).
\]
Thus coarse dependence means that at least one retained atom is already generated
by the remainder of the evidence body.
\textup{\textsc{Status: Lean-proved in
\texttt{PEL4/MatroidEvidenceCircuits.lean}.}}
\end{proposition}

\begin{theorem}[Complete circuit classification]
A body \(C\subseteq E\) is a channel circuit if and only if there exist distinct
\(x,y\in E\) such that
\[
  \chi(x)=\chi(y)
  \qquad\text{and}\qquad
  C=\{x,y\}.
\]
Equivalently, the circuits of the coarse 4-PEL partition matroid are exactly the
two-element same-polarity pairs.
\textup{\textsc{Status: Lean-proved in
\texttt{PEL4/MatroidEvidenceCircuits.lean}.}}
\end{theorem}
\begin{proof}[Proof sketch]
Every distinct same-polarity pair is dependent. Any dependent subset of that
pair must contain two distinct elements, so it must contain both members and
therefore equals the pair. Conversely, if \(C\) is minimally dependent, choose
two distinct same-polarity witnesses \(x,y\in C\). The pair \(\{x,y\}\) is
already dependent, hence minimality forces every member of \(C\) to lie in that
pair. Since both witnesses lie in \(C\), equality follows.
\end{proof}

This theorem turns the informal phrase ``redundant evidence'' into an exact
structural claim. At the coarse bilateral projection, a circuit does not mean
that one item of evidence is false, weak, or dispensable in every epistemic
sense. It means only that two distinct atoms occupy the same one-bit direction
seen by the FDE kernel.

\begin{proposition}[Circuit closure redundancy]
If \(x\neq y\) and \(\chi(x)=\chi(y)\), then for every \(e\in E\),
\[
  e\in\operatorname{cl}_{\chi}(\{x,y\})
  \quad\Longleftrightarrow\quad
  e\in\operatorname{cl}_{\chi}(\{x\}).
\]
By symmetry the same holds with \(\{y\}\). Hence removing either member of a
circuit leaves the coarse generated channel unchanged.
\textup{\textsc{Status: Lean-proved.}}
\end{proposition}

The circuit-elimination behaviour is correspondingly transparent. Suppose
\(\{x,y\}\) and \(\{x,z\}\) are distinct circuits sharing \(x\). Then all three
atoms have the same polarity; if \(y\neq z\), the pair \(\{y,z\}\) is itself a
circuit contained in the union after removing the common element. This is the
two-channel instance of the standard circuit-elimination pattern
\cite{Oxley2011}.

\begin{proposition}[Two-channel circuit elimination]
Let \(x,y,z\) be pairwise distinct with
\(\chi(x)=\chi(y)=\chi(z)\). Then
\[
  \{x,y\},\{x,z\}\text{ circuits}
  \quad\Longrightarrow\quad
  \{y,z\}\text{ is a circuit}.
\]
\textup{\textsc{Status: Lean-proved.}}
\end{proposition}

\subsection{Contraction as accepted background evidence}

Deletion and contraction are the canonical minor operations of matroid theory.
The current repository interface, however, presents only a closure operator and
has no explicit ground-set field. A full minor API would therefore be premature:
removing an atom from the carrier cannot yet be represented without changing the
ambient type or extending \texttt{MatroidClosure}. Gate~3 instead records the
part of contraction semantics that is expressible without pretending that this
missing structure already exists.

Fix an evidence atom \(c\) and treat it as accepted background evidence. Define
\[
  \operatorname{cl}_{\chi}^{/c}(A)
  :=\operatorname{cl}_{\chi}(A\cup\{c\}).
\]
This is a contraction-style background closure rather than a packaged matroid
minor.

\begin{proposition}[Parallel evidence becomes loop-like under background]
For any atom \(e\),
\[
  \chi(c)=\chi(e)
  \quad\Longrightarrow\quad
  e\in\operatorname{cl}_{\chi}^{/c}(\varnothing),
\]
whereas
\[
  \chi(c)\neq\chi(e)
  \quad\Longrightarrow\quad
  e\notin\operatorname{cl}_{\chi}^{/c}(\varnothing).
\]
Thus accepting one positive atom as background makes every positive parallel
atom generated without further evidence, while negative atoms remain a separate
information direction; and conversely for negative background evidence.
\textup{\textsc{Status: Lean-proved.}}
\end{proposition}

This is the strongest epistemic reading justified by the present formalism. In
a genuine contraction of a partition matroid, the parallel partners of the
contracted element become loops. Here the formal result captures exactly the
corresponding closure fact: once one channel representative is fixed in the
background, additional same-channel atoms contribute no new coarse rank.

\paragraph{Gate-3 boundary.}
The classification is exact only after the positive/negative projection. At a
finer level, two positive reports may differ in reliability, causal origin,
probability mass, inferential route, or independence of source. Such information
can split a coarse two-element circuit into an independent fine-grained set.
Consequently, ``minimal evidential redundancy'' is always indexed to a chosen
projection.

\paragraph{Next structural question.}
Gate~4 should promote the lightweight closure presentation to an explicit
finite-ground-set matroid layer, or bridge the project to a standard matroid API.
Only then can deletion and contraction be stated as actual minors rather than
background-closure probes. That gate would allow a precise comparison between
source removal, background acceptance, rank loss, loop creation, and preservation
of the four-valued state.
